\scanpage{56}
% L01-02
Pour que la \deleted{\unsure{bonne}} $\unsure{G|C}$-application \unsure{convienne} $\pi_E^{(i)}\circ\varphi_{\unsure{\mathscr{R}|\mathcal{R}|R}}$ passe au quotient \inserted{$\mathrm{Fac}_i(X)\ (=\underline{S},\underline{A},\underline{F})$}, il faut et s. que
% L03-05
\diagram{\begin{tikzcd}[ampersand replacement=\&]
\underline{\mathscr{R}} \arrow[r,"\varphi_{\mathscr{R}}"] \arrow[d,"\pi_{\underline{X}}^{(i)}"']
\& \mathrm{Drap}_{012}(E) \arrow[d,"\pi_E^{(i)}"] \\
\mathrm{Fac}_i(X) \arrow[r] \& \mathrm{Fac}_i(E)
\end{tikzcd}}
\[i=0,1,2\]

% L06-07
\deleted{$\unsure{(s_0,f_0)}$ \unsure{l'image}} le stabilisateur dans $G$ de $\unsure{f_0^{(i)}|\underline{f}_0^{(i)}}$ contienne celui de $f_0^{(i)}$, ce qui s'écrit
% L08-10
\[
\begin{aligned}
\deleted{\unsure{\sigma_0(s_0|\sigma_0(r_0}}\qquad
&\sigma_0(a_0)=a_0 &&\sigma_0(f_0)=f_0\\
\sigma_1(s_0)=s_0
&&&\sigma_1(f_0)=f_0\\
\sigma_2(s_0)=s_0\qquad
&\sigma_2(a_0)=a_0
\end{aligned}
\]
% archive sheet number: 55
